\section{Уточнённая постановка задачи}
\subsection{Назначение и исходные требования}
Разработать объектно-ориентированную программную систему, автоматизирующую ведение общего календаря фирмы и позволяющую исследовать актуализацию календаря при пошаговом изменении модельного времени. Исходная постановка требует визуализации с движущимися объектами, задания параметров перед экспериментом и статистического моделирования неопределённых факторов.

Фирма состоит из $K$ отделов, $5\le K\le9$. Во внутреннем календаре представлены ежедневные планёрки отделов, еженедельные совещания руководителей, командировки, тематические семинары и другие встречи. У мероприятия имеются название, тип, начало и окончание, участники, организатор, место и важность. Внутренние мероприятия используют помещения отделов или конференц-зал. Командировка не занимает внутреннее помещение.

Пользователь добавляет новые планы, уточняет существующие и получает сведения о событиях за день или период дня по фирме, отделу и сотруднику. Проверяется одновременное участие сотрудника в разных мероприятиях и пересечение бронирований помещения. Выявленный конфликт должен быть объяснён пользователю; желателен автоматический подбор нескольких способов его разрешения. Возможность изменения существующих планов зависит от категории пользователя.

Период опыта составляет $N$ календарных дней, $7\le N\le30$, шаг --- 30 или 60 минут. По завершении формируется отчёт о событиях фирмы за период. Уведомления участникам планируются по важности мероприятий.

\subsection{Принятые проектные решения}
В текущем проекте реализовано локальное приложение. Состав отделов, сотрудников и помещений постоянен на протяжении опыта. В каждом отделе первоначально руководителем является первый сотрудник. Отдельные кадровые и психологические процессы не входят в модель.

При включённом начальном календаре все рабочие дни периода заполняются 30-минутными планёрками каждого отдела. Они начинаются одновременно в начале рабочего дня, но проходят в разных помещениях и не конфликтуют. В первый рабочий день каждой календарной недели внутри периода после планёрок назначается 30-минутное совещание руководителей в конференц-зале. При выключенном пресете календарь первоначально пуст. Пользователь может дополнить его до продвижения времени.

Периодические семинары имеют отдельный шаблон и создаются во время прогона с указанной заблаговременностью. Их участники представляют минимум два отдела. Если дата нерабочая, она переносится вперёд. Случайные заявки и попытки изменения расписания генерируются независимо от визуальной скорости.

Модельное время хранится с точностью до минуты, со смещением UTC+03:00. Фактические часы компьютера не определяют ход событий. Начатые и завершённые мероприятия нельзя редактировать; завершённые исключаются из активного календаря, сохраняясь в архиве и истории.

\subsection{Роли и границы ручного управления}
До запуска выбирается роль оператора. В текущем интерфейсе оператор связан с сотрудником с внутренним ID 1; выбор произвольной учётной записи не реализован. Роль --- условие учебной модели, а не сетевая аутентификация.
\begin{itemize}
\item Секретарь может изменять весь календарь.
\item Руководитель может редактировать планы в разрешённом отделе. Проверяются организатор, участники и принадлежность помещения до и после изменения.
\item Сотрудник может создать мероприятие от своего имени, но не изменяет уже назначенные события.
\end{itemize}
Добавление, редактирование и выбор варианта переноса ставят прогон на паузу. Продолжение выполняется явно. Переводы сотрудников, назначения руководителей и отдельная кнопка отмены мероприятия в интерфейсе отсутствуют. Операция отмены сохранена в календарном сервисе для внутренних сценариев и проверок.

\subsection{Существенное ограничение относительно задания}
\textbf{Электронная почта имитируется.} Программа рассчитывает сроки и адресатов, сохраняет тему, текст, момент отправки и доставки, но не отправляет сообщения через SMTP. Модель считает доставку мгновенной и успешной. Поэтому требование настоящей электронной рассылки полностью не выполнено; применение имитации необходимо согласовать с преподавателем либо дополнить канал реальной отправкой. В настоящем отчёте эти действия не называются реальной доставкой.

\subsection{Исходные материалы}
Постановка уточнена по двум предоставленным документам: «Практикум. Моделирование.
Общая постановка задачи» и «Система автоматизации функций секретаря».
Сведения о реализации сверены с заголовками и исходниками текущего проекта,
а результаты исследования --- с полными записями прогонов calendar-v7.
